\begin{proposition}[二点纤维的中心型跳跃与非中心二值性的秩缺口]\label{prop:window6-foldbin-gauge-center-vs-charge-separation}
对分辨率 \(m\ge 1\)，令
\(
\mathcal{G}^{\mathrm{bin}}_m:=\Aut(\Fold^{\mathrm{bin}}_m)
\)
为二进制区间折叠诱导的 fold-gauge 规范群，并记纤维多重度
\(
d^{\mathrm{bin}}_m(w):=\card{(\Fold^{\mathrm{bin}}_m)^{-1}(w)}
\)（\(w\in X_m\)）。
定义计数
\[
s_m:=\#\{w\in X_m:\ d^{\mathrm{bin}}_m(w)\ge 2\},
\qquad
t_m:=\#\{w\in X_m:\ d^{\mathrm{bin}}_m(w)=2\}.
\]
则成立：
\begin{enumerate}
\normalfont
  \item \textbf{中心的完全分类。}规范群中心满足自然同构
  \[
  \boxed{\;Z(\mathcal{G}^{\mathrm{bin}}_m)\ \cong\ (\ZZ_2)^{t_m}.\;}
  \]
  其生成元可取为每条二点纤维上的唯一非平凡交换（单纤维 swap）；任意满足 \(d^{\mathrm{bin}}_m(w)\ge 3\) 的纤维不贡献非平凡中心元。
  \item \textbf{电荷格与“可实现跳跃”的严格分离。}由定理 \ref{thm:window6-foldbin-gauge-abelianization-even-parity} 有
  \[
  \boxed{\;\mathrm{Hom}(\mathcal{G}^{\mathrm{bin}}_m,\ZZ_2)\ \cong\ (\ZZ_2)^{s_m}.\;}
  \]
  并且对任意中心元 \(z\in Z(\mathcal{G}^{\mathrm{bin}}_m)\) 与任意二值电荷 \(\psi\in\mathrm{Hom}(\mathcal{G}^{\mathrm{bin}}_m,\ZZ_2)\)，取值 \(\psi(z)\in\ZZ_2\) 给出一条规范单射
  \[
  \boxed{\;Z(\mathcal{G}^{\mathrm{bin}}_m)\hookrightarrow \mathrm{Hom}(\mathcal{G}^{\mathrm{bin}}_m,\ZZ_2)^\vee.\;}
  \]
  在上述同构下，其像恰为“仅支撑在二点纤维坐标上的子空间”。
  因此商空间 \((\ZZ_2)^{s_m}/(\ZZ_2)^{t_m}\) 具有秩 \(s_m-t_m\)，并给出一个内禀的“非中心二值性”指标：它刻画了可作为电荷存在但不能由任何中心型二值跳跃算子在不引入外加标签的口径下实现的二值方向。
\end{enumerate}
\leanverified{binFold\_boundary\_count\_m6}
\leanverified{binFold\_boundary\_count\_m7}
\leanverified{binFold\_min\_mult\_m6}
\leanverified{window6\_six\_s\_count}
\leanverified{window6\_t\_count}
\leanverified{window6\_noncentral\_binary\_rank\_gap}
\leanverified{window6\_rank\_gap\_certificate}
\leanverified{paper\_window6\_rank\_gap\_extended}
\leanverified{paper\_gu\_window6\_boundary\_certificate}
\leanverified{paper\_window6\_foldbin\_gauge\_center\_vs\_charge\_separation}
\end{proposition}
\begin{proof}
由纤维分解，存在自然同构
\(
\mathcal{G}^{\mathrm{bin}}_m\cong \prod_{w\in X_m}\mathfrak S_{d^{\mathrm{bin}}_m(w)}
\)。
中心满足
\(
Z(\prod_w H_w)=\prod_w Z(H_w)
\)。
而 \(Z(\mathfrak S_d)\) 在 \(d\ge 3\) 时平凡，在 \(d=2\) 时同构于 \(\mathfrak S_2\cong\ZZ_2\)，在 \(d=1\) 时平凡；故中心同构于所有 \(d^{\mathrm{bin}}_m(w)=2\) 因子的直积，即 \((\ZZ_2)^{t_m}\)，并给出生成元描述。
\par
对电荷格：当 \(d\ge 2\) 时 \(\mathfrak S_d^{\mathrm{ab}}\cong\ZZ_2\)（由符号同态给出），当 \(d=1\) 时 Abel 化平凡；故
\(
\mathrm{Hom}(\mathcal{G}^{\mathrm{bin}}_m,\ZZ_2)\cong (\ZZ_2)^{s_m}
\)，与定理 \ref{thm:window6-foldbin-gauge-abelianization-even-parity} 一致。
\par
对单射：取值映射
\(
Z(\mathcal{G}^{\mathrm{bin}}_m)\times \mathrm{Hom}(\mathcal{G}^{\mathrm{bin}}_m,\ZZ_2)\to \ZZ_2
\)
双线性，故诱导同态
\(
Z(\mathcal{G}^{\mathrm{bin}}_m)\to \mathrm{Hom}(\mathcal{G}^{\mathrm{bin}}_m,\ZZ_2)^\vee
\)。
在直积分解下，中心仅来自 \(\mathfrak S_2\) 因子，而该因子上的唯一非平凡元在符号表示下取值为 \(1\)，从而像仅支撑在二点纤维坐标上，并且对每个非平凡中心元存在某个二点纤维坐标的符号电荷检测到其取值为 \(1\)，故该同态为单射。
秩缺口陈述由维数计数立得。证毕。
\end{proof}

\begin{remark}[window-\texorpdfstring{$6$}{6} 的秩缺口证书]\label{rem:window6-foldbin-noncentral-binary-rank-gap}
在 window-\(6\) 上，纤维尺寸直方图为 \(2:8,3:4,4:9\)（推论 \ref{cor:fold-bin-m6-fiber-hist}），故
\(
s_6=8+4+9=21
\)、
\(
t_6=8
\)，从而非中心二值性指标的秩为 \(s_6-t_6=13\)。
\end{remark}

\begin{theorem}[window-\texorpdfstring{$6$}{6} 的 fold-gauge 中心由八条二点纤维完全生成，且 boundary sheet parity 形成规范直和因子]\label{thm:window6-foldbin-gauge-center-boundary-direct-sum}
沿用推论 \ref{cor:fold-bin-m6-fiber-hist} 与推论 \ref{cor:fold-bin-boundary-center-z2} 的记号。记
\[
\mathcal G_6^{\mathrm{bin}}:=\Aut(\Fold_6^{\mathrm{bin}}),
\qquad
X_{6,2}^{\mathrm{int}}
:=
\{\,w\in X_6\setminus X_6^{\mathrm{bdry}}:\ d^{\mathrm{bin}}_6(w)=2\,\}.
\]
对每个满足 \(d^{\mathrm{bin}}_6(w)=2\) 的纤维，记其唯一非平凡交换为 \(\tau_w\)。再定义
\[
Z_{\partial}:=\langle \tau_w:\ w\in X_6^{\mathrm{bdry}}\rangle,
\qquad
Z_{\mathrm{int}}:=\langle \tau_w:\ w\in X_{6,2}^{\mathrm{int}}\rangle.
\]
则有自然同构
\[
\boxed{\;
Z(\mathcal G_6^{\mathrm{bin}})
=
Z_{\partial}\oplus Z_{\mathrm{int}}
\cong
(\ZZ_2)^3\oplus(\ZZ_2)^5
\cong
(\ZZ_2)^8.\;}
\]
并且 Abel 化电荷格分解为
\[
\boxed{\;
\bigl(\mathcal G_6^{\mathrm{bin}}\bigr)^{\mathrm{ab}}
\cong
Z_{\partial}\oplus Z_{\mathrm{int}}\oplus E_{\mathrm{sgn}},
\qquad
E_{\mathrm{sgn}}\cong (\ZZ_2)^{13}.\;}
\]
其中 \(E_{\mathrm{sgn}}\) 恰由全部 \(d^{\mathrm{bin}}_6(w)=3\) 与 \(d^{\mathrm{bin}}_6(w)=4\) 纤维的符号像生成。
\leanverified{paper\_window6\_foldbin\_gauge\_center\_boundary\_direct\_sum}
\end{theorem}
\begin{proof}
由推论 \ref{cor:fold-bin-m6-fiber-hist}，
window-\(6\) 的纤维尺寸直方图为
\[
2:8,\qquad 3:4,\qquad 4:9.
\]
因此
\[
\mathcal G_6^{\mathrm{bin}}
\cong
\mathfrak S_2^{\,8}\times \mathfrak S_3^{\,4}\times \mathfrak S_4^{\,9}.
\]
由命题 \ref{prop:window6-foldbin-gauge-center-vs-charge-separation} 的中心分类，
\[
Z(\mathcal G_6^{\mathrm{bin}})\cong (\ZZ_2)^{t_6}=(\ZZ_2)^8.
\]

另一方面，推论 \ref{cor:fold-bin-boundary-center-z2} 已给出
\[
Z_{\partial}\cong (\ZZ_2)^3.
\]
又由于全部二点纤维共有 \(8\) 条，而 boundary 纤维中恰有 \(3\) 条属于二点纤维，故
\[
\card{X_{6,2}^{\mathrm{int}}}=8-3=5,
\]
从而
\[
Z_{\mathrm{int}}\cong (\ZZ_2)^5.
\]
由于不同纤维对应于直积中的不同 \(\mathfrak S_2\) 因子，故边界二点纤维与内部二点纤维所生成的中心子群彼此独立，并且
\[
Z(\mathcal G_6^{\mathrm{bin}})
=
\prod_{w:\,d^{\mathrm{bin}}_6(w)=2} Z(\mathfrak S_2)
=
Z_{\partial}\oplus Z_{\mathrm{int}}.
\]
这就得到第一条分解。

再看 Abel 化。由定理 \ref{thm:window6-foldbin-gauge-abelianization-even-parity}，
\[
\bigl(\mathcal G_6^{\mathrm{bin}}\bigr)^{\mathrm{ab}}
\cong
(\ZZ_2)^{21}.
\]
更具体地，由直积分解与对称群 Abel 化，
\[
\bigl(\mathcal G_6^{\mathrm{bin}}\bigr)^{\mathrm{ab}}
\cong
\bigl(\mathfrak S_2^{\mathrm{ab}}\bigr)^{8}
\oplus
\bigl(\mathfrak S_3^{\mathrm{ab}}\bigr)^{4}
\oplus
\bigl(\mathfrak S_4^{\mathrm{ab}}\bigr)^{9}
\cong
(\ZZ_2)^8\oplus(\ZZ_2)^4\oplus(\ZZ_2)^9.
\]
其中前 \(8\) 个 \(\ZZ_2\) 因子正对应于全部二点纤维，因此在上式下恰可识别为
\[
Z_{\partial}\oplus Z_{\mathrm{int}}.
\]
余下
\[
(\ZZ_2)^4\oplus(\ZZ_2)^9\cong (\ZZ_2)^{13}
\]
正是所有三点与四点纤维的符号像，记为 \(E_{\mathrm{sgn}}\)。故
\[
\bigl(\mathcal G_6^{\mathrm{bin}}\bigr)^{\mathrm{ab}}
\cong
Z_{\partial}\oplus Z_{\mathrm{int}}\oplus E_{\mathrm{sgn}}.
\]
证毕。
\end{proof}

\endinput
